\documentclass[10pt]{article}
\usepackage[T1]{fontenc}
\usepackage[utf8]{inputenc}
\usepackage[french]{babel}
\usepackage{amsmath,amssymb,mathrsfs,mathpazo}
\usepackage{geometry}
\geometry{a4paper,margin=2cm}
\usepackage{graphicx,array,tabularx,tabularray,colortbl,multirow,multicol,booktabs,enumitem,pifont,ulem,fancybox,diagbox,cancel}
\usepackage{pstricks,pstricks-add,pst-plot,pst-tree,pst-node,pst-text,pst-func,pst-math,pst-bspline,pst-3dplot,pst-eucl,pst-solides3d}
\let\pscalculate\undefined
\usepackage{pst-bezier}
\usepackage{tikz,pgfplots,tkz-tab}
\usetikzlibrary{arrows,arrows.meta,patterns,calc,positioning,trees,decorations.markings,shapes.geometric,shapes.misc,shapes.arrows,intersections,decorations.pathmorphing,decorations.pathreplacing,backgrounds,shadows,babel}
\pgfplotsset{compat=1.18}
\usepgfplotslibrary{fillbetween}
\usepackage[np]{numprint}
\usepackage{listings,makecell,pas-tableur,twoopt}
\let\barre\undefined
\usepackage[pstricks]{tablvar}
\definecolor{darkgreen}{RGB}{0,100,0}
\providecommand{\red}{\color{red}}
\providecommand{\blue}{\color{blue}}
\providecommand{\green}{\color{green}}
\providecommand{\black}{\color{black}}
\providecommand{\white}{\color{white}}
\providecommand{\gray}{\color{gray}}
\providecommand{\cyan}{\color{cyan}}
\providecommand{\magenta}{\color{magenta}}
\providecommand{\yellow}{\color{yellow}}
\providecommand{\violet}{\color{violet}}
\providecommand{\orange}{\color{orange}}
\providecommand{\euro}{\texteuro}
\providecommand{\EUR}{\texteuro}
\providecommand{\ieme}{\textsuperscript{e}}
\providecommand{\degres}{\ensuremath{^\circ}}
\providecommand{\degree}{\ensuremath{^\circ}}
\providecommand{\vv}[1]{\overrightarrow{#1}}
\providecommand{\Coord}[1]{#1}
\providecommand{\decosix}{}
\providecommand{\starredbullet}{\ensuremath{\bullet}}
\pagestyle{empty}
\setlength{\parindent}{0pt}
\usepackage{adjustbox}\DeclareUnicodeCharacter{2236}{:}\DeclareUnicodeCharacter{2190}{\ensuremath{\leftarrow}}
\providecommand{\R}{\mathbb{R}}
\providecommand{\N}{\mathbb{N}}
\providecommand{\Z}{\mathbb{Z}}
\providecommand{\Q}{\mathbb{Q}}
\providecommand{\D}{\mathbb{D}}
\providecommand{\C}{\mathbb{C}}
\providecommand{\e}{{\mathrm{e}}}
\providecommand{\vect}[1]{\overrightarrow{#1}}
\providecommand{\vectt}[1]{\overrightarrow{\mathrm{#1}}}
\providecommand{\vvt}[1]{\overrightarrow{\mathrm{#1}}}
\providecommand{\np}[2][]{{#2~#1}}
\providecommand{\vv}[1]{{\overrightarrow{#1}}}
\renewcommand{\d}{{\mathrm{d}}}
\providecommand{\degres}{{^\circ}}
\providecommand{\degre}{{^\circ}}
\providecommand{\blue}{\color{blue}}
\providecommand{\red}{\color{red}}
\providecommand{\enskip}{\quad}
\providecommand{\llbracket}{\lbrack\!\lbrack}
\providecommand{\rrbracket}{\rbrack\!\rbrack}
\providecommand{\euro}{{\text{€}}}
\providecommand{\no}{{\text{n°}}}
\providecommand{\interval}[2]{{[#1;#2]}}
\providecommand{\intervFO}[2]{{[#1;#2[}}
\providecommand{\dotfill}{\ldots}

\providecommand{\ds}{\displaystyle}
\providecommand{\textding}[1]{#1}
\providecommand{\Oij}{\(\mathrm{O};\vec{\imath},\vec{\jmath}\)}
\providecommand{\Oijk}{\(\mathrm{O};\vec{\imath},\vec{\jmath},\vec{k}\)}
\providecommand{\Ouv}{\(\mathrm{O};\vec{u},\vec{v}\)}
\providecommand{\intervFF}[2]{\ensuremath{[#1;#2]}}
\renewcommand{\euro}{\ifmmode\text{\texteuro}\else\texteuro\fi}
\DeclareUnicodeCharacter{00B0}{\ifmmode^{\circ}\else\textdegree\fi}
\DeclareUnicodeCharacter{20AC}{\ifmmode\text{\texteuro}\else\texteuro\fi}
% Énoncés : Éducation nationale. Transcription et corrigés : APMEP, Denis Vergès et contributeurs. Découpage et classement : CoopMaths, Sébastien Lozano, Éric Elter et contributeurs. Adaptation de la présentation : Maths Explorer.
% CC BY-SA : https://coopmaths.fr/www/about/
% https://coopmaths.fr/alea/static/bac/2024/tex/bac_2024_11_sujet2_amsud_4.tex
\begin{document}
\noindent\textit{APMEP / CoopMaths. CC BY-SA. Présentation adaptée par Maths Explorer.}\par\medskip
\medskip

L'espace est muni d'un repère orthonormé \Oijk.

On considère les trois points A(3~;~0~;~0), B(0~;~2~;~0) et C(0~;~0~;~2).

\begin{center}
\psset{unit=0.8cm,arrowsize=2pt 3}
\begin{pspicture}(-2.5,-1.8)(3,2.5)
\psline[linewidth=1.25pt]{->}(1,0)
\psline[linewidth=1.25pt]{->}(0,1)
\psline[linewidth=1.25pt]{->}(-0.6,-0.6)
\uput[u](0.5,0){\footnotesize $\vect{\jmath}$}
\uput[l](0,0.5){\footnotesize $\vect{k}$}
\uput[ul](-0.3,-0.3){\footnotesize $\vect{\imath}$}
\psline[linewidth=1.25pt]{->}(1,0)
\psline(3,0)\psline(0,2.5)\psline(-2,-2)
\uput[ul](-1.8,-1.8){A}\uput[u](2,0){B}\uput[l](0,2){C}
\psdots(-1.8,-1.8)(2,0)(0,2)\uput[dr](0,0){O}
\end{pspicture}
\end{center}

L'objectif de cet exercice est de démontrer la propriété suivante :

\og Le carré de l'aire du triangle ABC est égal à la somme des carrés des aires des trois autres faces du tétraèdre OABC \fg,

\bigskip

\textbf{Partie 1 : Distance du point O au plan (ABC)}

\medskip

\begin{enumerate}
\item Démontrer que le vecteur $\vect{n}(2~;~3~;~3)$ est normal au plan (ABC).
\item Démontrer qu'une équation cartésienne du plan (ABC) est : $2x + 3y + 3z - 6 = 0$.
\item Donner une représentation paramétrique de la droite $d$ passant par O et de vecteur directeur $\vect{n}$.
\item On note H le point d'intersection de la droite $d$ et du plan (ABC).

Déterminer les coordonnées du point H.
\item En déduire que la distance du point O au plan (ABC) est égale à $\dfrac{3\sqrt{22}}{11}$ 
\end{enumerate}

\bigskip

\textbf{Partie 2 : Démonstration de la propriété}

\medskip

\begin{enumerate}
\item Démontrer que le volume du tétraèdre OABC est égal à 2.
\item En déduire que l'aire du triangle ABC est égale à $\sqrt{22}$.
\item Démontrer que pour le tétraèdre OABC, \og le carré de l'aire du triangle ABC est égal à la somme des carrés des aires des trois autres faces du tétraèdre \fg.

On rappelle que le volume d'un tétraèdre est donné par $V = \dfrac13B \times h$ où $B$ est l'aire d'une 
base du tétraèdre et $h$ est la hauteur relative à cette base.
\end{enumerate}

\end{document}
